%% APA 7th Edition — Student Pitch (no title page, no abstract)
%% Compile: pdflatex -> biber -> pdflatex -> pdflatex
\documentclass[12pt]{article}

\usepackage[margin=1in]{geometry}
\usepackage{setspace}
\doublespacing

\usepackage[american]{babel}
\usepackage{csquotes}
\usepackage[style=apa, backend=biber]{biblatex}
\DeclareLanguageMapping{american}{american-apa}
\addbibresource{references.bib}

\usepackage{hyperref}
\hypersetup{colorlinks=true, linkcolor=black, citecolor=black, urlcolor=blue}

\usepackage{titlesec}
% APA Level 1: centered, bold
\titleformat{\section}[block]{\normalfont\bfseries\centering}{}{0em}{}
% APA Level 2: left-aligned, bold
\titleformat{\subsection}[block]{\normalfont\bfseries}{}{0em}{}
% APA Level 3: left-aligned, bold italic, inline
\titleformat{\subsubsection}[runin]{\normalfont\bfseries\itshape}{}{0em}{}[.]

\titlespacing*{\section}{0pt}{12pt}{6pt}
\titlespacing*{\subsection}{0pt}{12pt}{6pt}
\titlespacing*{\subsubsection}{0pt}{6pt}{6pt}

\usepackage{parskip}
\setlength{\parskip}{0pt}
\setlength{\parindent}{0.5in}

\pagestyle{plain}

% ── Document ────────────────────────────────────────────────────────────────
\begin{document}

% Simple APA-style title block (no separate title page)
\begin{center}
  {\bfseries Large Lecture Model: A Behaviorally Grounded Course Intelligence
  Platform for Student Self-Regulation}\\[6pt]
  Tae Wook (Terry) Kim\\
  Massachusetts Institute of Technology\\
  CMS.594: Education Technology Studio
\end{center}

\vspace{6pt}

% ============================================================
\section{Introduction}
% ============================================================

Course material for a single large-lecture course is scattered across Canvas,
Gradescope, Panopto, and Piazza---each with its own interface---so a student
preparing for an exam cannot easily cross-reference which lectures introduced a
concept, which problem sets assessed it, and which discussion threads clarified
it. The \textit{Large Lecture Model} (LLM) addresses this by aggregating all
course content into a unified knowledge base exposed through two interfaces:
(1) an AI chatbot that answers natural-language questions with cited, grounded
responses, and (2) an assignment management dashboard linking deadlines to
topics and lecture materials. The prototype targets MIT 6.1220 (Design and
Analysis of Algorithms) and CMS.594 (Education Technology Studio) to validate
the system across a large STEM lecture and a smaller HASS course; the
architecture is parameterized to generalize to any course.

% ============================================================
\section{Concept and Final Product Vision}
% ============================================================

The system has three layers. \textbf{Data ingestion} scrapers pull content from
all four platforms through a pipeline of chunking, PII anonymization via spaCy
named-entity recognition, embedding generation, and LLM-based topic tagging.
The \textbf{knowledge layer} pairs ChromaDB semantic search (filtered by content
type and an academic integrity flag) with a Neo4j graph that connects Topics,
Lectures, Assignments, and Discussion Threads via edges (\texttt{COVERS},
\texttt{ASSESSED\_BY}, \texttt{PREREQ\_OF}); a hybrid retriever combines both.
\textbf{User interfaces} consist of a Next.js chatbot backed by a FastAPI
streaming server and an assignment dashboard with a topic map; a guardrail layer
blocks solution requests while permitting conceptual questions.

The current prototype has Canvas ingestion, vector retrieval, and a CLI
chatbot. The final product will add Panopto and Piazza ingestion with full
anonymization; the Neo4j graph builder and hybrid retriever; a tuned guardrail
with instructor-configurable protection windows (\texttt{courses.yaml}); and a
Next.js web interface with streaming citations, inline feedback, and reusable
packaging for any instructor to deploy.

% ============================================================
\section{Theoretical Grounding}
% ============================================================

\textbf{Cognitive load.} Platform fragmentation imposes extraneous cognitive
load \parencite{sweller1988}---students spend working memory navigating
interfaces rather than learning. A unified natural-language interface removes
this friction \parencite{mayer2009}, freeing attention for the material itself.

\textbf{Self-regulated learning and retrieval practice.} \textcite{zimmerman2002}
describes self-regulated learners as those who set goals, monitor progress, and
adapt strategy; the chatbot and topic-map dashboard scaffold exactly this cycle.
\textcite{roediger2006} showed active recall produces markedly better long-term
retention than re-reading---asking the chatbot is retrieval practice embedded in
a real study context. The graph's prerequisite edges further support the
connected knowledge structures that underlie expert transfer \parencite{chi1981}.

\textbf{Behavioral concept design.} A core design principle is drawn from
course readings \parencite{cms594staff2025}: the system models behaviors---
Posting, Asking, Submitting, Grading---in an event-based log. Because
``behaviors hold still even when platforms do not''
\parencite[p.~1]{cms594staff2025}, the same log serves the chatbot,
dashboard, and future tools without re-scraping or schema migration.

% ============================================================
\section{Design Tasks}
% ============================================================

The following tasks are needed to reach a finalized product, ordered by
dependency:

\begin{enumerate}
  \item Complete Panopto and Piazza scrapers with integrated PII anonymization;
        build the Neo4j graph from LLM-extracted syllabus relationships.
  \item Implement the hybrid vector+graph retriever behind a common interface.
  \item Tune the academic integrity guardrail with per-assignment, per-date-window
        sensitivity configured via \texttt{courses.yaml}.
  \item Build the Next.js web interface: streaming chatbot with citations,
        assignment dashboard, per-response feedback, and staff monitoring view;
        audit all UI components against WCAG 2.1 AA.
  \item Write setup documentation and configuration guide for any instructor.
\end{enumerate}

% ============================================================
\section{Stakeholder Interview Plan}
% ============================================================

\subsection{Students (Primary Users)}

Students currently enrolled in 6.1220 and CMS.594---the direct users---will be
asked how they currently find relevant materials and what would make them trust
or distrust a chatbot's answer. Feedback sought: navigation pain points, trust
signals, and academic integrity concerns to calibrate the guardrail.

\subsection{Course Staff (Instructors and TAs)}

The instructor of record and one or two TAs from each target course. Their
buy-in is a prerequisite for deployment. Key questions: What consumes the most
TA office hours? Which AI assistance is acceptable? Are there content types
(exam rubrics, partial solutions) that should never reach the system? Feedback
sought: content scope restrictions, guardrail calibration, and instructor-facing
features worth the setup cost.

\subsection{Students with Disabilities / Accessibility Consultants}

Students who use assistive technologies or MIT Digital Accessibility Services
staff, engaged during design rather than as a retrofit. Feedback sought:
specific accessibility barriers and modalities (voice input, simplified output)
addressable before the final build.

% ============================================================
\section{Accessibility, Viewpoint, Localization, and Feedback}
% ============================================================

\textbf{Accessibility.} The web interface will meet WCAG 2.1 AA: semantic HTML
for screen readers, full keyboard navigability, sufficient color contrast, and
careful focus management for streaming responses. A plain-text output mode will
be available for users with cognitive or attention-related needs.

\textbf{Viewpoint, localization, and feedback.} The event-log architecture is
viewpoint-neutral: a personal study view, aggregate coverage summary, and
future advising view all project the same log without privileging one use
\parencite{cms594staff2025}. The prototype is English-only, but the LLM layer
handles multilingual queries natively. Per-response rating and flagging are
logged without PII and surfaced to course staff as an aggregate monitoring view.

% ============================================================
\section{References}
\printbibliography[heading=none]
% ============================================================

\end{document}
